\clearpage
\begin{figure*}[t]
\centering
\begin{tcolorbox}[
  colframe=black,
  colback=gray!10,
  arc=2mm,
  boxrule=1.5pt,
  title=\textbf{Prompts},
  fonttitle=\bfseries,
]
\small
You are a deliverable planning assistant. Two agents collaborate to achieve a shared task goal. Your role is to turn the task goal into a concrete ``deliverable blueprint'': what the final output must CONTAIN (sections, fields, entities, feature lists, parameters, \ldots), so that the agents can actually produce it. You are generating content planning items that make the implicit deliverable contents explicit. 

The goal of this step is to FORCE specificity: 

- NOT ``include all relevant X''

- BUT ``include X1, X2, X3 \ldots''

\par\smallskip

\textbf{[Input]}

\textbf{Task Scenario:}

- This contains the scenario description that sets the context for the collaboration between the two agents.

- This contains the task goal that the two agents are collaborating to achieve.

- This contains the profile of the two agents.
\par\smallskip

\textbf{[Planning Guidelines]}

\textbf{1. Each item MUST specify concrete deliverable contents (not evaluation language).}

- Treat each item as a ``what must be written/encoded'' instruction.

- (Example) If the goal is a document: items should describe the document outline (sections) and the exact contents each section must include.

- (Example) If the goal is a configuration/spec file: items should describe the file schema (sections/keys) and the exact domain-specific entries that must appear.
\par\smallskip

\textbf{2. Force explicit domain content (invent plausible specifics when the goal implies them).}

- Do NOT say ``all relevant \ldots'' or ``all required \ldots''.

- DO provide the explicit list of feature groups, features, fields, parameters, \ldots\ the output should contain (as a planned baseline).

- Use the scenario domain to choose realistic, specific items.

- Avoid placeholders like ``TBD'', ``as appropriate'', ``etc.'', ``and more''.
\par\smallskip

\textbf{3. Each item MUST be atomic and simple.}

- One item should test exactly one requirement.

- Do NOT bundle multiple conditions into one item.

- Each item should be clear and concise; the content should consist only of core, significant elements.

- The content MUST NOT contain trivial elements.

- Minimize the amount of item contents.
\par\smallskip

\textbf{4. You must generate EXACTLY \$MAX\_REQUIREMENTS\$ requirements.}

- Even if more are possible, include only the \$MAX\_REQUIREMENTS\$ most important ones.
\par\smallskip
\textbf{[Output Format]}

Return a JSON object with EXACTLY the following structure (no extra text, no markdown formatting):

\begin{verbatim}
{
  "requirements": [
    "<planning item>",
    "...",
    ...
  ]
}
\end{verbatim}

\par\smallskip


\textbf{[Input]}

\textbf{Task Scenario:}

\$SCENARIO\_JSON\$

\par\smallskip


\textbf{Output:}


\end{tcolorbox}
\caption{A prompt used for generating agent profiles and task goals.}
\label{fig:prompt_dataset_2}
\end{figure*}